\section{Diskussion und Grenzen}

\subsection{Technische Machbarkeit und Robustheit}

Die Hauptfrage kann für den kontrollierten Alurohrdatensatz und den
nachgewiesenen Matrix-/Smoke-/Replaypfad grundsätzlich positiv beantwortet
werden. Die Verarbeitungskette kann aus einer Bildsequenz Masken und Kameras
erzeugen, objektbezogene Gaussians trainieren, ein Mesh ableiten und daraus
eine lokale Centerline samt B-Spline exportieren. Der wissenschaftlich
relevante Befund ist dabei nicht nur ein einzelner erfolgreicher Lauf. Die im
Matrix-/Replaypfad wirksamen Quality-Gates für Masken, Kameramodelle, ideale
Bilddomäne, Eval-Split und Ergebnisarchive verhindern mehrere zuvor
beobachtete stille oder verspätete Fehler; der Inline-Vollauf muss vor der
Abgabe noch auf dieselbe Gate-/Warp-Kette gebracht oder separat begrenzt werden.

Robustheit bedeutet in dieser Arbeit nicht Fehlerfreiheit unter jeder
denkbaren Eingabe. Sie bedeutet, dass bekannte Voraussetzungen geprüft,
unvollständige Routen eindeutig markiert, teure Stufen an kontrollierten
Grenzen wiederholt und Ergebnisse reproduzierbar archiviert werden. Die
Untersuchung nur eines Testvideos begrenzt weiterhin die externe Validität.

\subsection{Masken, Domänen und Geometrie}

Die zentrale Domänentrennung löst einen methodischen Widerspruch: SfM profitiert
von unmaskierten Hintergrundmerkmalen, während STS und objektmaskierte
Metriken eine semantisch begrenzte ideale Bilddomäne benötigen. Die gemeinsame
Transformation von Bildern, Kameras und Masken stellt diese Konsistenz her.

Masken sind dennoch keine geometrische Referenz. Sie gewichten oder selektieren
Bildpixel und Objektzuordnungen. Position, Skalierung und Orientierung der
Gaussians, Kamerafehler, Surface-Sampling und Poisson-Rekonstruktion können die
3D-Geometrie weiterhin verändern. Deshalb ist die zunächst plausible Aussage,
eine niedrige Auflösung genüge bei korrekten Masken und registriertem COLMAP,
nur als technische Hypothese zulässig. Für einen erfolgreichen Lauf müssen alle
Datenverträge gelten; für geometrische Genauigkeit wäre zusätzlich eine
unabhängige Referenz notwendig.

\subsection{Kameramodell, FPS und Auflösung}

Die Matrix zeigt, dass Kameramodell, zeitliche Samplingdichte und Auflösung
nicht isoliert interpretiert werden dürfen. Die 2-FPS-Profile enthalten weniger
zeitlich benachbarte Ansichten; bei dünnen Objekten kann dies die SAM-
Propagation und die verfügbare Überdeckung erschweren. Die Beobachtung ist
jedoch nicht monoton: Einzelne Low- und 720p-Läufe bleiben maskenstabil.

Die nativen Bildmetriken werden in unterschiedlichen Pixelräumen berechnet.
Downsampling kann Bildfehler und Maskengrenzen mitteln. Ein strengerer
Auflösungsvergleich müsste alle Renderings auf eine gemeinsame Zielauflösung
und identische Eval-Views bringen. Die derzeitige Matrix ist deshalb ein
Engineering-Screening.

Die Auflösungen dürfen daher nicht zu einer globalen Qualitätsrangliste
zusammengefasst werden. Ein Low-Lauf kann durch geglättete Kanten und kleinere
hochfrequente Fehler höhere PSNR-/SSIM-Werte oder niedrigere LPIPS-Werte
erreichen, ohne eine bessere 3D-Oberfläche zu besitzen. Wissenschaftlich
belastbarer sind gepaarte Vergleiche innerhalb derselben Auflösung,
Kameramodell- und FPS-Konfiguration. Die neue Boxplot-Auswertung ergänzt die
Mittelwerte um die Streuung der einzelnen Testansichten; sie ersetzt aber
nicht die geometrische Referenzmessung.

Die Bildtafeln des Pipeline-Durchlaufs präzisieren diese Einschränkung. Trotz
des gleichen SIFT-Merkmalslimits sind in der Punktwolken-Übersicht kaum
Auflösungsunterschiede sichtbar; lokal weist die \texttt{720p}-Punktwolke in
der linken unteren Wandecke jedoch mehr sichtbare Punkte auf. Eine räumlich
breitere Verteilung der bei begrenztem Merkmalsbudget gefundenen homologen
Punkte könnte die Track- und SfM-Stabilität verbessern. Dies bleibt eine
plausible Hypothese zur internen Geometriekonsistenz und ist keine Aussage
über die absolute Richtigkeit der Rekonstruktion. Ebenso widerspricht ein
solcher möglicher SfM-Vorteil höheren SSIM-Werten einer niedrigeren Auflösung
nicht: SSIM bewertet die Ähnlichkeit der gerenderten Bildansicht im jeweiligen
Pixelraum, nicht die räumliche Verteilung der SfM-Tracks oder die geometrische
Richtigkeit des späteren Meshes. Die sichtbaren Floater-, Rand- und
Centerline-Unterschiede zeigen deshalb, warum Bildmetrik und Geometrie getrennt
bewertet werden müssen.

Der direkte PINHOLE-Lauf darf nicht als Nachweis einer Entzerrung interpretiert
werden. Er erzwingt auf den Rohbildern die Annahme einer idealen Kamera. Der
Produktionsweg mit \texttt{OPENCV} (Standard seit der Matrixauswertung,
Abschnitt~\ref{sec:produktionskonfiguration}) schätzt dagegen zunächst eine
Verzeichnung und erzeugt anschließend eine separate ideale Szene; dasselbe
gilt für die historische \texttt{SIMPLE\_RADIAL}-Baseline. OPENCV kann
zusätzliche Verzeichnungsanteile modellieren, besitzt aber mehr
Freiheitsgrade. Aus den Bildmetriken allein lässt sich daher keine generell
beste Kamerakonfiguration ableiten; der gewählte Standard ist eine
datengestützte Projektentscheidung für 720p bei 5 FPS und bleibt für
veränderte Datensätze neu zu prüfen. Die von COLMAP ausgegebenen
Verzeichnungskoeffizienten sind aus den Bildkorrespondenzen geschätzte
Parameter, nicht unabhängig gemessene Objektivdaten. Ein niedriger mittlerer
Reprojektionsfehler belegt daher die interne Geometriekonsistenz, aber nicht
die physikalische Richtigkeit der Entzerrung. Für einen Vermessungsnachweis
ist eine externe Kamerakalibrierung mit Referenzmustern oder bekannten
Kalibrierparametern erforderlich; diese liegt in der Projektarbeit nicht vor.

\subsection{Original-GS und SuGaR-Coarse}

Die Vierfeld-Ablation zeigt, dass nicht nur die verwendete Tiefenroute, sondern
auch die zusätzliche SuGaR-Coarse-Optimierung die resultierende Oberfläche
verändert. Route A bewahrt den STS-Zustand am direktesten und ist für das Ziel
einer stabilen Centerline deshalb der plausibelste Produktionsstandard. Mehr
Vertices oder sichtbar feinere Strukturen sind nicht automatisch besser, wenn
sie zugleich Spitzen oder veränderte Randstützen erzeugen.

Die Entscheidung ist bewusst vorläufig. Die Meshdistanzen vergleichen nur
Varianten untereinander. Sie geben keinen Abstand zum realen Alurohr an. Die
zwischenzeitlich abgeschlossene zwölfteilige SuGaR-Coarse-Folgematrix zeigt,
dass der reparierte Coarse-Renderpfad über beide Kameramodelle, drei
Auflösungen und zwei FPS-Profile technisch durchläuft. Sie liefert damit eine
breitere technische Evidenz als der einzelne Smoke-Checkpoint. Da in keinem
der zwölf Läufe ein `refined.ply` exportiert wurde, ist die Aussage weiterhin
auf SuGaR-Coarse beschränkt; eine SuGaR-Refined-Auswertung bleibt offen.

Der gepaarte Delta-Vergleich der zwölf Follow-up-Konfigurationen ist deshalb
als Rendervergleich zu lesen: SuGaR-Coarse wird jeweils mit dem STS-Splat
derselben Konfiguration verglichen. Die beobachteten Unterschiede können neben
der zusätzlichen SuGaR-Optimierung auch aus dem unterschiedlichen Renderpfad
entstehen. Die Delta-Grafik beweist somit weder eine allgemeine Überlegenheit
noch eine Unterlegenheit von SuGaR und keine Meshgenauigkeit.

\subsection{Autopilot, Replay und Batchfähigkeit}

Der Matrixrunner belegt die serielle Batchverarbeitung deklarierter
Konfigurationen. Der Replay reduziert Wiederholungsaufwand, verlangt aber
vollständige und kompatible Eingabeartefakte an der gewählten Schrittgrenze.
Der Autopilot füllt interaktive Entscheidungen mit Standardwerten und kann den
manuellen GCP-Breakpoint überspringen. Diese drei Funktionen sind verwandt,
aber nicht identisch.

Vor einer endgültigen Aussage zur nutzerseitigen Autonomie muss ein
Autopilot-Vollauf oder ein gleichwertiger Vertragstest archiviert werden. Die
Containerisierung und nichtinteraktive Matrix bereiten einen Serverbetrieb vor;
ein Deployment, Mehrbenutzerbetrieb und Betriebsmonitoring wurden nicht
evaluiert. Benutzerfreundlichkeit wird deshalb nur als reduzierte Zahl
manueller Entscheidungen, nicht als formale Usability-Messung bewertet.

\subsection{Laufzeit und Effizienz}

Der Zeit-/Qualitätsquotient ist ebenfalls sekundär. Er kombiniert normalisierte
PSNR-, SSIM- und invertierte LPIPS-Werte mit gleichen Gewichten und verwendet
nur die archivierten STS-bis-Postprocess-Zeiten. SAM3- und COLMAP-Zeit sowie
Container-Overheads der Rendering- und Metrikschritte sind nicht vollständig
enthalten. Vollständige End-to-End-Zeiten liegen zwischenzeitlich aus dem
Verifikationsbatch \path{matrix_e2e_verifikation_260826} vor
(Abschnitt~\ref{sec:laufzeiten-e2e}); sie ergänzen die Ressourcenplanung,
ändern aber nichts an der Segmentbasis des Quotienten. Für die
Bachelorarbeit müssen Rohmetriken und geometrische
Referenzmetriken höher gewichtet werden.

Die Effizienzanalyse ist trotzdem praktisch relevant: Ein automatischer
Screeninglauf kann ungeeignete Kombinationen aussortieren, bevor eine
hochaufgelöste Produktionsrechnung erfolgt. Wegen unvollständiger
End-to-End-Laufzeiten bleibt die Zeit-/Qualitätsmatrix im Anhang und dient nicht
als primäre wissenschaftliche Rangliste.

\subsection{Centerline, Georeferenzierung und externer Geltungsbereich}

Die größte technische Restaufgabe ist die echte Georeferenzierung. Ein
Translation-Fallback prüft zwar die Ausführung, sagt aber nichts über Lage-
oder Höhengenauigkeit aus. Ebenso sind die Testdaten keine reale Rohr- oder
Kabeltrasse und enthalten keine unabhängige GNSS-Referenzkurve.

Auch die lokale Centerline ist bisher nur als erzeugtes Artefakt, nicht als
geometrisch richtige Scheitelachse validiert. Der \texttt{single}-Modus ist für
eine einzelne Trasse robust, verwirft aber bewusst Netzwerktopologie. Eine
Übertragung auf Abzweigungen erfordert neue Referenzdaten und eigene
Graphmetriken.

Für die Bachelorarbeit folgen daraus vier priorisierte Schritte: reale
Aufnahmebedingungen, echte GCP-/GNSS-Daten, validierte relative
4x4-Transformation sowie RMSE- und Hausdorff-Distanz einer extrahierten
Centerline gegen eine unabhängige Referenzkurve. Erst damit kann die
ingenieurtechnische Toleranz geprüft werden.

\subsection{Beantwortung der Forschungsfragen}

\begin{enumerate}
  \item Die konsistente Übergabe erfordert ein erhaltenes Originalmodell, eine
	  getrennte ideale Bilddomäne, denselben Warp für Masken, einen
	  nichtleeren gematchten Eval-Split und explizite Statusprüfungen aller
	  Ergebnisstufen.
  \item Kameramodell, FPS und Auflösung verändern Pipelineerfolg und
	  Ansichtsmetriken, jedoch nicht monoton. Die aktuelle Matrix erlaubt ein
	  Screening, aber keinen isolierten Kausal- oder Geometrienachweis.
  \item Route A bewahrt die Original-STS-Gaussians direkter und zeigt im
	  relativen Vergleich weniger unerwünschte Oberflächenveränderung. Ihre
	  absolute geometrische Richtigkeit bleibt offen.
  \item Replay und Matrixrunner sind für reproduzierbare Wiederholungen und
	  Batchläufe nachgewiesen. Der Autopilot ist implementiert; sein
	  vollständiger End-to-End-Nachweis bleibt vor Abgabe zu archivieren.
\end{enumerate}
